(1) 镜面沿一臂移动 $\Delta L$，往返传播距离变化 $2\Delta L$，对应真空等效光程差变化 $\Delta(\mathrm{OPD})=2n\Delta L$。因光束同一介质传播，额外光程需乘相应折射率；仅计算单程 $\Delta L$ 会漏掉干涉仪中反射来回两次的距离。
(2) 相位差为 $\Delta\phi=(2\pi/\lambda_0)\Delta(\mathrm{OPD})=4\pi n\Delta L/\lambda_0$。每当相位差增加 $2\pi$，干涉亮暗条纹经历一个完整周期，因此扫过的完整条纹数 $N=|\Delta\phi|/(2\pi)=2n|\Delta L|/\lambda_0$；末端若非整数周期，实际计数还需记录剩余相位。
(3) 若镜面以速度 $v$ 移动且 $|v|$ 足够小以忽略动态传播高阶修正，时间相位变化率为 $\dot{\Delta\phi}=4\pi nv/\lambda_0$，所以条纹扫过频率 $f_{\rm beat}=|\dot{\Delta\phi}|/(2\pi)=2n|v|/\lambda_0$。在真空 $n=1$，这与低速反射镜造成的近似双程多普勒频移 $|\Delta f|\simeq2|v|/\lambda_0$ 一致。
(4) 若光源光谱有有限宽度，探测强度干涉项由不同频率分量的复相干叠加决定。典型相干长度 $L_c$ 约与中心波长平方除以谱宽成正比，若两臂光程差 $\gg L_c$，不同频率相位平均相互抵消，条纹可见度趋低；单频公式只给每个频率自身的相位，不保证多频混合仍形成高对比度干涉。
